\vspace*{-1cm}
\section{Gestión del proyecto}
\label{sec:gestion-proyecto}

\subsection{Modelo de ciclo de vida}
Se deberá indicar la metodología de desarrollo seguida, es decir, si se ha llevado a cabo un modelo de ciclo de vida en cascada, incremental, etc.

\subsection{Planificación}
Se identificarán las tareas a realizar, su descripción, sus fechas estimadas de inicio y de fin, sus precedencias, los recursos necesarios, estimación de esfuerzo, etc. Se podrán utilizar diagramas de Gantt, PERT / CPM, etc. Deberán incluirse (si procede) también los diferentes hitos a lo largo del proyecto. Para que el estudiante tenga la experiencia real de planificar y controlar su cumplimiento, es \textbf{importante} que la planificación se realice con anterioridad a la finalización del proyecto.

Se pueden encontrar explicaciones tanto sobre los diagramas que se indican en esta sección como los que se mencionan en los siguientes apartados en la \href{https://administracionelectronica.gob.es/pae_Home/dam/jcr:da7d91fa-d6bd-467c-be32-a72e27c603b3/METRICA_V3_Tecnicas.pdf}{\color{blue}{documentación sobre técnicas de la metodología Métrica 3}}.

\subsection{Presupuesto (si procede)}
En caso de que el desarrollo del TFG en sí mismo implique costes, o tenga sentido un trabajo con proyección a futuro, se indicará qué personal (con sus perfiles, dedicación, etc.), licencias de software, infraestructuras, etc. serán necesarios, así como sus precios y costes correspondientes.